\section{Open Questions}

The five open questions below are genuinely unresolved by this replication;
each is accompanied by the basis on which it is open and concrete next steps.

\begin{enumerate}

\item \textbf{Does an independently regenerated reversible circuit for the elliptic curve controlled point-addition (Roetteler et al.\ Algorithm~1) yield the same leading Toffoli coefficient $224 = 4\cdot 32 + 2\cdot 16 + 4\cdot 16$, and does its regression against $n$ produce a subleading term compatible with the paper's $+2045 n^2$?}

\emph{Basis.} This replication takes Table~1's per-primitive counts on the
authors' authority and reconstructs the composition. The subleading $+2045 n^2$
term is fit-derived from LIQUi$|\rangle$ runs we cannot re-execute (F\# source
not released). A primitive-regenerative check would isolate whether the
$+2045 n^2$ constant is intrinsic to Algorithm~1's bookkeeping or an artifact
of specific LIQUi$|\rangle$ optimizations.

\emph{Next steps.} Implement the controlled EC point-addition circuit in
Qualtran (using its \texttt{ModInverse}, \texttt{ModMul}, \texttt{ModSqu}
bloqs), decompose to Toffoli$+$AND, count for $n \in \{110, 160, 192, 224, 256,
384, 521\}$, regress $a n^2 \log n + b n^2$, and compare $a$ to 224 and $b$ to
2045. Prior open-source Qualtran ECC bloqs implement Litinski~2023's windowed
variant, not Roetteler~2017's un-windowed variant, so a new bloq class
\texttt{RoettelerECAdd} would need to be authored ($\sim 500$--$1000$ lines Python).

\item \textbf{How does the Roetteler~2017 $(448 \log_2 n + 4090) n^3$ Toffoli cost compare, in a common primitive-cost framework, against modern windowed-arithmetic ECDLP-Shor variants (Litinski~2023, H\"aner-Roetteler-Svore, Gidney's windowed Montgomery), and against Ekera-H\aa stad phase-estimation-reduction techniques applied to ECDLP?}

\emph{Basis.} Our Qualtran cross-check shows Litinski~2023 gives $\sim 170\times$
smaller Toffoli count for $n = 256$, but the two counts are produced by
different tools, different algorithms, and different windowing conventions ---
the ratio is suggestive not definitive. Ekera-H\aa stad-style multiple-runs-with-
reduced-QPE has been analyzed for RSA/DLP but not systematically for ECDLP
against Roetteler's baseline. H\"aner-Roetteler-Svore arithmetic variants may
change constants noticeably.

\emph{Next steps.} (1) Implement each candidate algorithm's controlled
point-addition in the same framework (Qualtran) with the same window-size
parameterization; (2) tabulate Toffoli count and qubit count at $n \in \{256,
384, 521\}$; (3) also count non-Clifford AND-computations separately;
(4) publish a side-by-side table with matched assumptions. Estimated effort:
2--3 person-weeks per algorithm.

\item \textbf{What is the physical resource cost of Roetteler~2017's ECDLP-Shor circuit under a specific fault-tolerant code family --- rotated surface code with lattice surgery, or newer LDPC codes (bicycle, hypergraph-product) --- for physical error rates $p \in \{10^{-3}, 10^{-4}\}$ and target logical failure per Toffoli of $10^{-15}$?}

\emph{Basis.} The paper (and this replication) report only \emph{logical}
Toffoli counts. Physical-cost translation multiplies logical Toffoli count by
(a) magic-state distillation factor ($10^2$--$10^4$ physical Toffolis per
logical, depending on distillation depth), (b) lattice-surgery routing
overhead, and (c) code-distance-squared for physical qubits. No published
Roetteler-plus-FT-overhead calculation exists at ECC-521 with 2024/2025
distillation protocols (e.g.\ Litinski's cultivation).

\emph{Next steps.} (1) Feed the Roetteler Toffoli count into Qualtran's
\texttt{PhysicalCost} estimator with a modern distillation config
(e.g.\ 15-to-1 with cultivation, code distance $\sim 25$); (2) also try qLDPC
estimators from Bravyi et al.\ 2024; (3) report physical qubits and wall-clock
time for a 1~MHz surface-code cycle. Expected outputs are of order $10^7$
physical qubits and hours-to-days runtime; the exercise is to pin down the
coefficient.

\item \textbf{Does ternary (base-3) reversible arithmetic --- as implemented for factoring in some qutrit proposals --- reduce Roetteler's per-point-addition Toffoli count for ECDLP over $\mathbb{F}_p$ when $p$ is not a power of 2? If not, does mixed-radix arithmetic (base 2 for storage, base 4 or higher for carry lookahead) help?}

\emph{Basis.} Roetteler 2017 uses pure binary Montgomery arithmetic. Ternary
reversible arithmetic has been explored for Shor's factoring (Bocharov et al.,
Anwar et al.) and can give asymptotic savings on some subroutines but pays a
cost in the $T$-gate cost of the qutrit magic state. The trade-off for the
modular-inversion-dominated ECDLP cost profile is not settled.

\emph{Next steps.} (1) Locate published ternary reversible mod-inverse circuit
constructions and their Toffoli-equivalent (or $T$-count) analyses;
(2) if available, evaluate their cost at $n = 256/384/521$ and compare to
Roetteler's Toffoli / equivalent-$T$; (3) if not available, implement a base-3
extended-binary-GCD analog in Qualtran-qutrit or write it out symbolically.
Estimated effort: 1 person-month; a null result (ternary does not help ECDLP)
would itself be publishable.

\item \textbf{How does Roetteler's ECDLP-Shor cost translate to isogeny-based cryptosystems (CSIDH, SQIsign) --- which are post-quantum candidates specifically because ordinary Shor does not directly attack them --- and to the isogeny-adapted Shor's variant (Kuperberg's / Regev's dihedral-hidden-subgroup approach)? At what $n$-analog do the two crossover in cost?}

\emph{Basis.} CSIDH and SQIsign resist standard Shor because their hard problem
(isogeny walk) is a hidden shift / abelian HSP not a discrete-log HSP.
Kuperberg's algorithm attacks them in subexponential time. Comparing the
actual Toffoli / qubit cost of Kuperberg-on-CSIDH vs.\ Roetteler-on-ECDLP at
equivalent classical-security levels quantifies how much post-quantum runway
CSIDH really provides.

\emph{Next steps.} (1) Enumerate the isogeny-graph diameter and step-cost for
CSIDH-512 / CSIDH-1024 / SQIsign; (2) plug into Kuperberg / Regev cost formulas
(see Bonnetain-Schrottenloher~2020, Peikert~2020); (3) compare Toffoli count at
128-bit and 256-bit classical security against Roetteler's ECDLP-Shor at NIST
P-256/P-521; (4) publish crossover. This would sharpen NIST's post-quantum
standardization confidence.

\end{enumerate}
